% ============================================================
%  SECTION 2 — Introduction to Cryptography
% ============================================================
\section{Introduction to Cryptography}

Cryptography is the study of mathematical techniques for securing communication
and data in the presence of adversaries. Its primary goal is to ensure the
confidentiality, integrity, authenticity, and non-repudiation of information.
We are mainly dealing with two subjects in it: Encryption \& Decryption.

\begin{smjdefinition}
In cryptography, \textbf{encryption} (more specifically, \textbf{encoding})
is the process of transforming information in a way that, ideally, only
authorized parties can decode.
\end{smjdefinition}

\subsection{RSA}

The RSA cryptosystem was introduced in 1977 by Rivest, Shamir, and Adleman
\cite{rsa1978}, from whom it takes its name. It is a public-key encryption
scheme designed to securely encode messages over insecure communication
channels by relying on the computational difficulty of factoring large
composite integers. Let's discuss how it works and what is the algorithm
behind it.

\subsection{Algorithm of RSA}

Take Alice and Bob. Alice and Bob both have two big prime numbers (in general
around 313 digits). Let's name the prime numbers of Alice $P_1$ and $P_2$,
and Bob's $P_3$ and $P_4$.

Both Alice and Bob are going to multiply their two big prime numbers to get a
new big semiprime number. Let's call them respectively $SP_1$ and $SP_2$.

\begin{smjdefinition}
A \textbf{symmetric} (private-key) cryptosystem uses a single secret key
shared between sender and receiver --- the same key encrypts and decrypts the
message. The problem is that both parties must first agree on the key, which
requires a secure channel to begin with.

An \textbf{asymmetric} (public-key) cryptosystem solves this by using a pair
of mathematically linked keys: a \textbf{public key}, shared openly and used
to encrypt, and a \textbf{private key}, kept secret and used to decrypt. The
two keys are linked such that encrypting with one can only be undone by the
other. Crucially, deriving the private key from the public key is
computationally infeasible.
\end{smjdefinition}

$P_1$ and $P_2$ are then called the private keys of Alice, $SP_1$ is her
public key. And similarly for Bob.

If Alice wants to send something to Bob, she is going to hash her message in
a way such that she uses his public key $SP_2$, and the only way to decrypt
her message, must be that her recipient knows the two big prime numbers $P_3$
and $P_4$, that created the public key $SP_2$. In our case, Bob knows them,
but no one else does. Both the public keys of Alice and Bob can be viewed by
anyone, as well as her hashed message. However, due to the fact that only Bob
knows what $P_3$ and $P_4$ are, he is the only one able to decode the hashed
message.

\begin{smjnote}
That we are not covering how the message is hashed, and we won't. This part
is not relevant for the rest of the article and thus is out of scope.
\end{smjnote}

\begin{smjnote}
RSA is an asymmetric key system because two different keys are used to send
and to receive the message.
\end{smjnote}

\subsection{How to Make RSA Irrelevant?}

The whole RSA is standing on one assumption: finding the prime factors of a
semiprime number is a hard problem. More precisely, no polynomial-time
classical algorithm for integer factorization has ever been found, and the
best known classical algorithms still run in sub-exponential but
super-polynomial time. For a $313$-digit semiprime number, even the most
powerful supercomputers would take thousands of years to factor it. Thus, to
break RSA, it is enough to find an efficient factoring algorithm --- one that
runs in polynomial time. So far, no one has been able to do so on a classical
computer.

And this is where Shor's algorithm comes in.

\begin{smjnote}
Many have probably heard that factoring a semiprime number is an NP problem.
While it is true, it does not really say anything about RSA being strong or
not. Without entering into too much detail, an NP problem is just a problem
where one can verify if a solution is correct in polynomial time \cite{arora2009},
which is indeed possible in prime factorization: one can check if the product
of 2 prime numbers is equal to the semiprime number. But it does not say
anything about whether a problem is in P or not --- any problem in P (problems
which we proved are solvable in polynomial time) is also in NP; summing
$3 + 5$ is also an NP problem.
\end{smjnote}

\subsection{Introduction to Shor's Algorithm}

Shor's algorithm, developed by Peter Shor in 1994 \cite{shor1994}, is a
quantum algorithm for efficiently factoring large integers. It runs in
polynomial time on a quantum computer and poses a fundamental threat to
classical cryptographic systems such as RSA, whose security relies on the
hardness of integer factorization.

But before explaining how it works, you may have noticed that there is a quite
scary word in its definition, ``quantum''. Thus, before explaining the
algorithm and how it is implemented, for the sake of the reader's comfort,
I will briefly introduce the world of Quantum Computing.